\chapter{Συμπεράσματα}

Η παρούσα μελέτη εξέτασε την αναγνώριση ηλεκτρομηχανικών ρυθμών στο σύστημα IEEE 39-Bus από ambient αποκρίσεις, με χρήση του αλγορίθμου N4SID και επακόλουθη συσταδοποίηση των screened εκτιμήσεων. Τα δύο order sweeps οδήγησαν σε παρόμοιες συγκεντρώσεις εκτιμήσεων στο επίπεδο $(\sigma,f)$, γεγονός που υποστηρίζει τη σταθερότητα των βασικών ευρημάτων ως προς την επιλεγμένη τάξη του μοντέλου. Η συχνότητα αρκετών ρυθμών αναφοράς του \textit{PowerFactory} ανακτήθηκε με ικανοποιητική ακρίβεια, ενώ η εκτίμηση της απόσβεσης παρουσίασε αρκετά μεγαλύτερη διασπορά.

Η συγκριτική εφαρμογή των αλγορίθμων έδειξε ότι δεν υπάρχει μία μέθοδος που υπερέχει σε όλα τα κριτήρια. Οι \textit{k-Means} και \textit{k-Medoids} αξιοποίησαν το σύνολο των εκτιμήσεων και παρείχαν σταθερές, εύκολα ερμηνεύσιμες λύσεις. Ο GMM παρείχε λεπτομερέστερο διαχωρισμό, με την τελική κάλυψη να εξαρτάται από την εφαρμογή του φίλτρου silhouette. Ο AHC οδήγησε σε δύο έως έξι τελικές συστάδες, διατηρώντας σχεδόν το σύνολο των εκτιμήσεων. Οι DBSCAN και OPTICS απέδωσαν συνεκτικές συστάδες, χαρακτηρίζοντας όμως σημαντικό ποσοστό των εκτιμήσεων ως θόρυβο. Ο HDBSCAN παρουσίασε πιο ισορροπημένη συμπεριφορά μεταξύ συνοχής, κάλυψης και απόρριψης θορύβου.

Οι χάρτες συσταδοποίησης και οι ποσοτικές μετρικές επιβεβαιώνουν ότι οι συστάδες αναπαράγουν τη συχνότητα αρκετών ρυθμών αναφοράς. Παράλληλα, εμφανίζονται συστάδες σε πολύ χαμηλές και υψηλές συχνότητες, όπου δεν αντιστοιχεί κάποιος ρυθμός αναφοράς. Το εύρημα αυτό, σε συνδυασμό με τις σχετικά μεγάλες τιμές MAD σε ορισμένες αντιστοιχίσεις, δείχνει ότι η συχνότητα αναγνωρίζεται γενικά πιο αξιόπιστα από την απόσβεση και ότι η αξιολόγηση πρέπει να βασίζεται συνδυαστικά.

Τα συμπεράσματα αφορούν το συγκεκριμένο συνθετικό ambient σενάριο, τις επιλεγμένες μετρήσεις ρεύματος και τάσης και τα εξεταζόμενα εύρη υπερπαραμέτρων των αλγορίθμων συσταδοποίησης. Επόμενο βήμα είναι η επανάληψη της διαδικασίας για δίκτυα που έχουν εισχωρήσει ανανεώσιμες πηγές ενέργειας.
